%----------------------------------------------------------------------------
\chapter{Introduction}\label{ch:intro}
%----------------------------------------------------------------------------

This project started as a request from the ELTE Cybersecurity Lab. The lab had a growing number of repositories on GitHub, a research group with multiple ongoing projects, publications spread across different venues, and no single place where any of this lived together. What they wanted was a portal that could pull all of it under one roof: a website that showed what the lab was working on, who was on the team, what had been published, and how the public repositories fit into the bigger picture. The aim was for someone arriving at the portal to understand the lab the same way they would understand a course: not by scrolling through a wall of repository names, but by reading something that was organised, paced, and easy to follow.

That last part shaped the project more than anything else. The lab manager described it clearly during the first meetings. Each repository should feel like a chapter in a book. Not a raw README dumped onto a page, but a structured walkthrough with a clear opening, a set of topics underneath it, and an order the reader could follow. A visitor should be able to land on a project, see a short overview, and then drill into the parts they care about without losing context. That image, repository as chapter, is what the portal was built around.

%-----------------------------------------------------------
\section{Why a Dynamic Data-Driven Approach}\label{sec:intro-dynamic}

The first version of the CodeBases page took a single repository and rendered the entire README on one long page, formatted the way GitHub renders it. That worked for one repository. It stopped working the moment a second one was added. The page became a wall of text. Two different projects had two completely different shapes, because nobody had agreed on what a README should look like. There was no way to navigate inside a project, no way to compare projects, no consistent visual identity, and certainly no chapter-like feel.

The fix turned into the central decision of the project. Instead of pulling raw Markdown and dumping it on the page, the README itself was given a structure. A small convention was agreed on for how a repository's README should be written: a top-level heading for the project, second-level headings for topics, third-level headings for sub-topics. With that convention in place, a parser could read any README that followed it and split the content into navigable sections. The CodeBases page then stopped being a wall of text and started behaving like a course outline. A sidebar listed the topics, the main panel showed the currently selected topic, and the user moved through the project the way they would move through a chapter.

Once that idea worked for one page, it became the pattern for the whole portal. Every other page was rebuilt around the same principle: separate the data from the component that displayed it. The Projects page reads from a single TypeScript file that lists every project, its description, its sub-projects, and the codebases each project relates to. The Research page reads from a file that lists each research area with its title, colour, intro, and full description. The Team page reads from a file that lists members by category. The Publications page reads from a list of publication entries. The Operations page reads from a file that describes the four roles and the three guiding principles. Even the navigation bar, the footer text, the contact details, and the gallery photos all sit in data files. None of this content lives inside a component. Components only read data and display it.

This separation was not done for elegance. It was done because the lab is going to keep growing. There will be new members, new publications, new projects, new repositories. If updating the portal required editing component code every time, the portal would slowly fall behind the lab. By keeping everything in data files, anyone on the lab can open a single \texttt{.ts} file, change a name or add an entry, and the portal updates. The components do not need to be touched. The architecture is designed for tens or hundreds of entries to be added later without anything breaking.

%-----------------------------------------------------------
\section{What the Portal Includes}\label{sec:intro-includes}

The finished portal has nine main pages reachable from the navigation bar: Projects, Research, Operations, CodeBases, Publications, Team, Mini Apps, and Gallery, plus a Contact page reachable from the footer. The header is shared across every page and includes the lab title, the navigation bar, and a theme toggle that switches the entire portal between light and dark mode. The footer is also shared and includes the lab's address, links to the university and faculty, and the contact and GitHub buttons.

A few pages are worth flagging here because they shaped the architecture more than the others. The CodeBases page is the one that drove the structured-README approach. The Projects page uses an interactive node graph where each project is a node on a diamond and clicking a node opens that project's detail panel. The Research page uses an interactive honeycomb where each research area is a hexagonal cell and clicking a cell opens that area's full description. The Operations page uses a two-by-two grid where each quadrant is a project role and clicking a role opens its description. Each of these pages was a design decision in its own right, and each was built so that adding a new project, area, or role is a single new entry in a data file.

The Mini Apps page hosts a small set of interactive cybersecurity tools, currently Ancient Cipher for classical encryption and decryption, and Security Gauntlet for a set of cybersecurity-themed mini-games. The intent is for this section to grow over time as more teaching material is added.

Beyond the pages themselves, the portal also has a GitHub organisation page (the \texttt{.github} repository's profile README) that serves as a landing page for visitors who arrive through GitHub rather than the portal. That page mirrors the structure of the portal: a short introduction, the research areas, featured projects with technology badges, and links back to the main portal.

%-----------------------------------------------------------
\section{Structure of This Document}\label{sec:intro-structure}

The rest of this document is organised as follows. Chapter~\ref{ch:user} is the user documentation. It walks through the portal page by page, describing what a visitor sees, what each clickable element does, and which data file controls what appears on each page. It is written for two audiences at once: visitors who want to understand the portal, and lab members who want to edit it without touching the code.

Chapter~\ref{ch:dev} is the developer documentation. It opens with the technology stack and the architectural decisions behind it, then explains the design choices for each part of the portal and why alternative approaches were rejected. It covers the data flow, the parser used for the CodeBases page, the pluggable component registry used on the Operations page, the theming system, and the routing setup. It closes with a section on testing, since the testing for this project was done manually and is short enough to belong inside the developer chapter rather than a separate one.

Chapter~\ref{ch:challenges} is the challenges and improvements chapter. It describes what was difficult during the project, where the current implementation has clear gaps (the contact form has no backend, the GitHub Actions automation needs an access token to run, there is no admin dashboard), and what would be reasonable next steps for whoever picks the project up next.
